\chapter{Limits} Used in defining the important concepts in calculus such as continunity, derivative of a function, the definite integral of a function.
\section{What are Limits}
It's easy to brush off limits due to their inherent simplicity in solving for them, but it's good to know what is happening when you say a limit equals something. This means that as you zoom in on that one area it is \textit{approaching}, the number it equals, despite it potentially never reaching it, still is the answer. The limit is the number you inch closer and closer to while zooming in.\\

\subsection{Inituitive Definition}
The \textbf{limit} of a function describes the behavior of a function near a particular value of x.\\
\\
When a limit heads towards infinity due to a vertical asymptote this is said to be a \textbf{infinite limit}.\\

\section{One-Sided Limits}
When a limit has a '-' sign that means the limit as it approaches from the \textit{left}.\\
When a limit has a '+' sign that means the limit as it approaches from the \textit{right}.\\
\\

\section{Trigonometric Functions}
There are four important limit functions:\\
1. $\lim\limits_{x \to c} sinx = sinc$\\
2. $\lim\limits_{x \to c} cosx = cosc$\\
3. $\lim\limits_{x \to 0} \frac{sinx}{x} = 1$\\
4. $\lim\limits_{x \to 0} \frac{1-cosx}{x} = 0$

\section{Continunity}
Three conditions must be met for a point to considered continuous:\\
1.  f(c) exists\\
2. $\lim\limits_{x \to c} f(x)$\\
3. $\lim\limits_{x \to c} f(x) = f(c)$

\newpage
\begin{center}
	\begin{tikzpicture}
		\begin{axis}[standard,
			xtick={-7,...,6},
			ytick={-4,...,5},
			samples=1000,
			xlabel={$x$},
			ylabel={$y$},
			xmin=-7,xmax=6,
			ymin=-4,ymax=5,
			clip=false]

			\addplot[
				<->,
				thick,
				] {x+2} node[anchor=south] {c};

			\coordinate (P) at (axis cs: 2, 4); % Define a coordinate at the point (0, 2)
			\draw[thick,fill=white] (P) circle (3pt) node[right] {b}; % Draw a hollow circle centered at P
		
			\coordinate (Q) at (axis cs: 2, 1); % Define a coordinate at the point (0, 2)
			\draw[thick,fill=black] (Q) circle (3pt) node[right] {a}; % Draw a hollow circle centered at P
		
		\end{axis}
	\end{tikzpicture}
\end{center}

In the graph above we can say multiple things:\\
1. A limit exists at $f(b)$ which is 4.\\
2. The limit \textit{at} $f(b)$ is 1 aka 'a'.\\
3. This line has a disconnected continunity. \\
\\
A test would be to imagine a pencil going along $f(x)$ without having to be picked up -- if it can accomplish this then it is continuous.


